\documentclass[11pt]{article}
\usepackage[margin=0.68in]{geometry}
\usepackage{lmodern}
\usepackage{microtype}
\usepackage{enumitem}
\usepackage[table]{xcolor}
\usepackage{titlesec}
\setlength{\parindent}{0pt}
\setlength{\parskip}{0.26em}
\setlist[itemize]{leftmargin=1.3em,itemsep=0.06em,topsep=0.1em}
\titleformat{\section}{\large\bfseries}{}{0pt}{}
\titlespacing*{\section}{0pt}{0.55em}{0.2em}
\pagestyle{plain}
\newenvironment{playpassage}{\begin{quote}\small\setlength{\parskip}{0.05em}\setlength{\parindent}{0pt}}{\end{quote}}
\newcommand{\speaker}[1]{\textbf{#1:}}
\begin{document}
\begin{center}
{\Large\bfseries U1L15SP --- Text Supplement}\par
{\LARGE\bfseries Competing Claims about Manhood}\par
{\large Two scenes for close reading; the whole play remains available}
\end{center}
\fcolorbox{black}{gray!8}{\begin{minipage}{0.94\linewidth}\small
\textbf{How to use this supplement}

You may use \textbf{anything accurate from \textit{Macbeth}} in your paper. The passages below are lines worth focusing on, not limits on your evidence. Quote exact words when they are printed here. You may accurately paraphrase other events you remember from watching the play. Do not invent quotations from memory.
\end{minipage}}

\section*{Passage 1: Act 1, Scene 5 --- Lady Macbeth Prepares Herself}
\textbf{Context:} Lady Macbeth has read Macbeth's report of the witches' greetings. A messenger announces that King Duncan will arrive at her castle that night. Alone, she calls on spirits as she prepares for what she intends.
\begin{playpassage}
\speaker{Lady Macbeth} The raven himself is hoarse\par
That croaks the fatal entrance of Duncan\par
Under my battlements. Come, you spirits\par
That tend on mortal thoughts, unsex me here,\par
And fill me from the crown to the toe top-full\par
Of direst cruelty! make thick my blood;\par
Stop up the access and passage to remorse,\par
That no compunctious visitings of nature\par
Shake my fell purpose, nor keep peace between\par
The effect and it! Come to my woman's breasts,\par
And take my milk for gall, you murdering ministers,\par
Wherever in your sightless substances\par
You wait on nature's mischief! Come, thick night,\par
And pall thee in the dunnest smoke of hell,\par
That my keen knife see not the wound it makes,\par
Nor heaven peep through the blanket of the dark,\par
To cry ``Hold, hold!''
\end{playpassage}
\section*{Words in context}
\begin{itemize}
\item \textbf{mortal thoughts}: deadly intentions; \textbf{unsex}: strip away traits associated with her sex;
\item \textbf{compunctious}: producing guilt or remorse; \textbf{visitings of nature}: natural feelings of pity;
\item \textbf{fell}: fierce or deadly; \textbf{gall}: bitter fluid, contrasted with milk;
\item \textbf{sightless substances}: invisible forms; \textbf{pall}: cover;
\item \textbf{dunnest}: darkest gray-brown.
\end{itemize}

\newpage
\section*{Passage 2: Act 1, Scene 7 --- Macbeth and Lady Macbeth Dispute Manhood}
\textbf{Context:} Macbeth has decided, ``We will proceed no further in this business.'' Lady Macbeth challenges his courage and love. Macbeth answers by drawing a boundary around what a man should dare.
\begin{playpassage}
\speaker{Lady Macbeth} Was the hope drunk\par
Wherein you dress'd yourself? hath it slept since?\par
And wakes it now, to look so green and pale\par
At what it did so freely? From this time\par
Such I account thy love. Art thou afeard\par
To be the same in thine own act and valour\par
As thou art in desire? Wouldst thou have that\par
Which thou esteem'st the ornament of life,\par
And live a coward in thine own esteem,\par
Letting ``I dare not'' wait upon ``I would,''\par
Like the poor cat i' the adage?

\speaker{Macbeth} Prithee, peace:\par
I dare do all that may become a man;\par
Who dares do more is none.

\speaker{Lady Macbeth} What beast was't, then,\par
That made you break this enterprise to me?\par
When you durst do it, then you were a man;\par
And, to be more than what you were, you would\par
Be so much more the man.
\end{playpassage}
\section*{Words in context}
\begin{itemize}
\item \textbf{green and pale}: sickly or afraid; \textbf{afeard}: afraid;
\item \textbf{valour}: courage; \textbf{ornament of life}: the crown or highest prize desired;
\item \textbf{esteem}: judgment or opinion; \textbf{adage}: proverb;
\item \textbf{become a man}: be fitting or proper for a man; \textbf{enterprise}: plan.
\end{itemize}

\section*{The rest of the play is available to you}
You may test the competing meanings against later choices and consequences: Duncan's murder, Macbeth's rule and further violence, Lady Macbeth's attempt to steady Macbeth, the banquet, her sleepwalking, either character's response to guilt, or their deaths. These are optional reminders, not conclusions.

A passage can reveal how a character uses a word without settling which meaning is right. Explain the meaning you infer, the evidence that tests it, and what real disagreement remains after the word is clarified.
\end{document}
